%% FullCourse_10Lecs -- Lecture 7, Topic 7.4: Economic Applications / Textual Analysis
%% Source: ./archive_pre_split/Lec07_LLM.tex (section 8)
%% Pass B: layer in refinements from
%%         ../../MiniCourse_8hr/Slides/Module3_LLM_Text/M3_T4_Textual_Analysis.tex
%%         (FinBERT/CentralBankRoBERTa, ECB Word2Prices, validation checklist)

%% Shared preamble for the FullCourse_10Lecs topic decks.
%%
%% Each topic .tex \input's this file, then defines its own \title, \date
%% override (if any), \graphicspath, and \begin{document}.
%%
%% Reconstructed verbatim from the inlined copy preserved in
%% Lec10_Case_Studies/Lec10_T8_Case6_DeepRamsey.tex (lines 23-190), which is
%% explicitly annotated there as "from ../shared_preamble.tex, copied verbatim".
%% Base preamble matches MiniCourse_8hr/Slides/shared_preamble.tex; this version
%% additionally provides \sectiondivider, the conceptbox/cautionbox/examplebox/
%% takeawaybox tcolorboxes, and \compacttables used across the split decks.

\documentclass[10pt,english,aspectratio=169]{beamer}

\usepackage{ae,aecompl}
\usepackage[T1]{fontenc}
\usepackage{textcomp}
\usepackage[utf8]{inputenc}
\usepackage{babel}
\usepackage{datetime2}

\setcounter{secnumdepth}{3}
\setcounter{tocdepth}{3}

\usepackage{amsmath, amssymb, amsthm, graphicx}
\usepackage{booktabs, multirow}
\usepackage{tikz}
\usepackage{pgfplots}
\pgfplotsset{compat=1.16}
\usepackage[most]{tcolorbox}
\tcbuselibrary{skins, breakable, theorems}
\usepackage{listings}
\usepackage{xcolor}

\lstset{
  basicstyle=\ttfamily\small,
  keywordstyle=\color{blue!80!black}\bfseries,
  commentstyle=\color{green!50!black}\itshape,
  stringstyle=\color{orange!80!black},
  backgroundcolor=\color{gray!8},
  frame=single,
  rulecolor=\color{gray!40},
  breaklines=true,
  showstringspaces=false,
  numbers=none,
}

\lstdefinestyle{pythonstyle}{
    language=Python,
    basicstyle=\ttfamily\footnotesize,
    keywordstyle=\color{blue!80!black}\bfseries,
    commentstyle=\color{green!50!black}\itshape,
    stringstyle=\color{orange!80!black},
    frame=single,
    breaklines=true,
    showstringspaces=false,
    numbers=left,
    numbersep=5pt,
    numberstyle=\tiny\color{gray},
    backgroundcolor=\color{gray!8},
    rulecolor=\color{gray!40}
}

\ifx\hypersetup\undefined
  \AtBeginDocument{%
    \hypersetup{bookmarks=false,breaklinks=true,pdfborder={0 0 0},colorlinks=true,
      linkcolor=DarkText,urlcolor=RedTitle}%
  }
\else
  \hypersetup{bookmarks=false,breaklinks=true,pdfborder={0 0 0},colorlinks=true,
    linkcolor=DarkText,urlcolor=RedTitle}%
\fi

\makeatletter
\newcommand\makebeamertitle{%
  \begin{frame}[plain]
    \vspace*{1.0cm}
    \centering
    {\usebeamerfont{title}\usebeamercolor[fg]{title}\inserttitle\par}
    \vspace{1.0cm}
    {\usebeamerfont{author}\usebeamercolor[fg]{author}\insertauthor\par}
    \vspace{0.2cm}
    {\usebeamerfont{date}\usebeamercolor[fg]{date}\insertdate\par}
  \end{frame}}%
\AtBeginDocument{%
  \let\origtableofcontents=\tableofcontents
  \def\tableofcontents{\@ifnextchar[{\origtableofcontents}{\gobbletableofcontents}}
  \def\gobbletableofcontents#1{\origtableofcontents}%
}

\usetheme{metropolis}
\usefonttheme{professionalfonts}

\definecolor{LightGray}{RGB}{230,230,230}
\definecolor{RedTitle}{RGB}{0,0,200}
\definecolor{VeryLightGray}{RGB}{245,245,245}
\definecolor{DarkText}{RGB}{0,0,0}

\setbeamercolor{background canvas}{bg=white}
\setbeamercolor{title}{fg=RedTitle}
\setbeamercolor{author}{fg=DarkText}
\setbeamercolor{date}{fg=DarkText}
\setbeamercolor{frametitle}{bg=VeryLightGray,fg=DarkText}
\setbeamercolor{normal text}{fg=DarkText}
\setbeamerfont{title}{size=\large}

\setbeamersize{text margin left=5mm, text margin right=5mm}
\addtobeamertemplate{frametitle}{}{\vspace{-0.5em}}
\newcommand{\reducevspace}{\vspace{-0.5cm}}

% Shared section divider and box styles for split topic decks.
\newcommand{\sectiondivider}[2]{%
\begin{frame}[plain,noframenumbering]
\begin{center}
\vfill
{\Large\textbf{#1}\par}
\vspace{0.35cm}
{\normalsize #2\par}
\vfill
\end{center}
\end{frame}
}

\newtcolorbox{conceptbox}[1][]{
  colback=blue!3,
  colframe=RedTitle!55,
  boxrule=0.35mm,
  arc=1mm,
  left=2mm,
  right=2mm,
  top=1mm,
  bottom=1mm,
  #1
}
\newtcolorbox{cautionbox}[1][]{
  colback=red!3,
  colframe=red!50!black,
  boxrule=0.35mm,
  arc=1mm,
  left=2mm,
  right=2mm,
  top=1mm,
  bottom=1mm,
  #1
}
\newtcolorbox{examplebox}[1][]{
  colback=green!3,
  colframe=green!45!black,
  boxrule=0.35mm,
  arc=1mm,
  left=2mm,
  right=2mm,
  top=1mm,
  bottom=1mm,
  #1
}
\newtcolorbox{takeawaybox}[1][]{
  colback=VeryLightGray,
  colframe=RedTitle!55,
  boxrule=0.35mm,
  arc=1mm,
  left=2mm,
  right=2mm,
  top=1mm,
  bottom=1mm,
  #1
}
\newcommand{\compacttables}{\renewcommand{\arraystretch}{1.15}}

\setbeamertemplate{navigation symbols}{}
\setbeamertemplate{footline}{%
  \hfill%
  \usebeamercolor[fg]{page number in head/foot}%
  \usebeamerfont{page number in head/foot}%
  \insertframenumber\,/\,\inserttotalframenumber\kern1em\vskip2pt%
}

\makeatother

\author{Zhigang Feng}
% "date with time": compile timestamp so each build is identifiable.
\date{\today\ \DTMcurrenttime}


\graphicspath{{./pic/}{../../MiniCourse_8hr/Slides/Module3_LLM_Text/pic/}{../}}

\title[AI for Econ Research]{AI for Economic Research: Dynamic Models, Language, and Agents\\
Lecture 7.4: Economic Applications: Textual Analysis}

\begin{document}

\makebeamertitle

%=============================================================================
\begin{frame}{Topic 7.4 Agenda}
\begin{enumerate}
  \item \textbf{Economic Applications} --- FOMC, SEC filings, news
  \medskip
  \item \textbf{Domain-Specific Models} --- FinBERT, CentralBankRoBERTa
  \medskip
  \item \textbf{Case Studies} --- ECB Word2Prices, validation checklist
\end{enumerate}

\bigskip
\footnotesize\textit{Arc: T1a Why$+$Tokens $\to$ T1b Embeddings $\to$ T2a Attention $\to$ T2b Transformer $\to$ T3a Training $\to$ T3b Limits $\to$ T4 Applications.}
\end{frame}

%=============================================================================
\section{Economic Applications}

\begin{frame}{Three Featured Case Studies}
\begin{enumerate}
\item \textbf{ECB Word2Prices} (Araujo, Bokan, Comazzi \& Lenza, 2025)
   \begin{itemize}
       \item Word2Vec on ECB statements + BVAR for inflation forecasting
   \end{itemize}

\medskip
\item \textbf{Patents and Worker Exposure} (Kogan et al., 2019)
   \begin{itemize}
       \item TF-IDF + cosine similarity between patents and occupational tasks
   \end{itemize}

\medskip
\item \textbf{Central Bank Sentiment with BERT} (Pfeifer \& Marohl, 2023; Ganguli et al., 2024)
   \begin{itemize}
       \item Fine-tuned RoBERTa for sentiment toward different economic agents
       \item PaECTER patent embeddings for innovation dynamics
   \end{itemize}
\end{enumerate}

\medskip
After the three deep dives, we briefly tour the wider ``application zoo'' and the validation pitfalls econometricians need to keep in mind.
\end{frame}

\begin{frame}{Case 1 --- ECB Word2Prices: Setup}
\textbf{Araujo, Bokan, Comazzi \& Lenza (2025), ECB WP No.\ 3047}

\medskip
\begin{itemize}
\item \textbf{Question:} can embeddings of ECB communications improve euro-area inflation forecasting?

\medskip
\item \textbf{Data:}
\begin{itemize}
    \item ECB press conference introductory statements, 2002Q1--2023Q4
    \item Target: HICPx (core inflation), monthly $\rightarrow$ quarterly
\end{itemize}

\medskip
\item \textbf{Method:}
\begin{itemize}
    \item For each quarter $q$, retrain Word2Vec on the corpus available \emph{up to} $q$ (no look-ahead)
    \item Average word vectors per statement; aggregate to a quarterly text vector
    \item Stack text vector with macro variables; fit a Bayesian VAR
    \item Forecast 1--4 quarters ahead
\end{itemize}

\medskip
\item \textbf{Why this design?} Re-training quarterly avoids the look-ahead bias that plagues most LLM-based forecast exercises, while remaining computationally cheap.
\end{itemize}
\end{frame}

\begin{frame}[fragile]{Case 1 --- Python Skeleton}
\begin{lstlisting}[language=Python, basicstyle=\scriptsize\ttfamily]
# Quarterly Word2Vec retraining for ECB inflation forecasting
from gensim.models import Word2Vec
from statsmodels.tsa.api import VAR
import numpy as np

forecasts = []
for q in quarters:
    # 1. Build corpus available at quarter q (no look-ahead)
    corpus_q = [s for s in all_statements if s.date <= q]
    # 2. Train Word2Vec on the available corpus
    w2v = Word2Vec(corpus_q, vector_size=100, window=5,
                   sg=0, min_count=2, epochs=20)
    # 3. Average word vectors per statement, then per quarter
    stmt_vec = [np.mean([w2v.wv[w] for w in s if w in w2v.wv], axis=0)
                for s in corpus_q]
    text_q = aggregate_quarterly(stmt_vec, dates)
    # 4. Stack with macro variables and fit a BVAR
    Y = np.column_stack([macro_q, text_q])
    bvar = VAR(Y).fit(maxlags=4)
    # 5. Forecast 1-4 quarters ahead
    forecasts.append(bvar.forecast(Y, steps=4))
\end{lstlisting}
\end{frame}

\begin{frame}{Case 1 --- Results and Lessons}
\begin{itemize}
\item \textbf{Forecast performance:}
\begin{itemize}
    \item BVAR with Word2Vec text features outperforms a baseline AR / BVAR without text at horizons of 1--4 quarters
    \item Sentiment dictionaries (e.g.\ Loughran-McDonald) deliver only marginal gains relative to the no-text baseline
    \item Heavy LLMs (BERT, fine-tuned) can match or beat Word2Vec but are much more expensive and harder to retrain quarterly --- look-ahead bias becomes a real risk
\end{itemize}

\medskip
\item \textbf{Take-aways for empirical economists:}
\begin{itemize}
    \item ``Lighter is better'' if it lets you respect the timing structure of the forecasting exercise
    \item Frequently-retrained Word2Vec is a strong baseline before reaching for an LLM
    \item Always document: training corpus cutoff at every step, vector dimension, window, retraining frequency
\end{itemize}

\medskip
\item Replicable framework with zero GPU requirements --- ideal for grad-student empirical work.
\end{itemize}
\end{frame}

\begin{frame}{Case 2 --- Patent Text and Worker Exposure}
\textbf{Kogan, Papanikolaou, Schmidt, Seegmiller (2019, 2023)}

\medskip
\begin{itemize}
\item \textbf{Question:} how does exposure to technological innovation affect workers' earnings and employment?

\medskip
\item \textbf{Method:}
\begin{itemize}
    \item Compute TF-IDF vectors for U.S.\ patent texts and for O*NET occupational task descriptions
    \item Measure cosine similarity between each occupation and each patent
    \item Aggregate up: occupation-level technological exposure
\end{itemize}

\medskip
\item \textbf{Findings:}
\begin{itemize}
    \item High-exposure occupations experience earnings declines and reduced employment stability
    \item Effects are present for both low- and high-wage workers
\end{itemize}

\medskip
\item \textbf{Why this is in our lecture:} a textbook example of how a \emph{simple} method (TF-IDF + cosine) on a \emph{thoughtfully chosen} corpus can answer a substantive economic question. You don't need transformers to make a contribution, but you do need a clean identification strategy.
\end{itemize}
\end{frame}

\begin{frame}{Why Generic LLMs Underperform on Financial Text}
\textbf{The vocabulary mismatch.} Generic BERT/GPT was trained on Wikipedia, web text, and books --- not on disclosure language.

\medskip
\begin{itemize}
  \item \textbf{Polarity flips.} ``Liability,'' ``aggressive,'' ``exposure'' are negative in everyday English but neutral or positive in accounting/finance. Loughran \& McDonald (2011) is the canonical correction of the older Harvard-IV-4 list.
  \medskip
  \item \textbf{Hedging language.} ``May,'' ``could,'' ``intend to'' are stop-words in generic NLP but encode legal disclosure stance in 10-K text.
  \medskip
  \item \textbf{Numbers in context.} Subword tokenizers split \texttt{1.2345} into digit tokens; without finance fine-tuning the model treats the magnitude as just another vocabulary item.
\end{itemize}

\medskip
\textbf{Two responses:} (1) train a domain-specific encoder (next slide), or (2) zero-shot prompt a frontier LLM and validate carefully (Hansen \& Kazinnik, 2023).
\end{frame}

\begin{frame}{Case 3a --- FinBERT and CentralBankRoBERTa: Two Domain Encoders}
\begin{center}
\renewcommand{\arraystretch}{1.25}
\footnotesize
\begin{tabular}{|p{2.6cm}|p{4.0cm}|p{6.4cm}|}
\hline
& \textbf{FinBERT} (Araci, 2019; Yang et al., 2020) & \textbf{CentralBankRoBERTa} (Pfeifer \& Marohl, 2023) \\ \hline
Backbone & BERT-base (12 layers, 110M params) & RoBERTa-base, fine-tuned \\ \hline
Pre-train corpus & Reuters TRC2-Financial + 10-K corpora & RoBERTa-base fine-tuned on $\sim$12k hand-labeled sentences (plus ECB/BIS pseudo-labels) \\ \hline
Task it ships for & 3-way sentiment (positive/negative/neutral) on financial sentences & Agent-specific sentiment: stance toward five macro agents (households, firms, financial sector, government, central bank) \\ \hline
Headline result & Beats LSTM/dictionary baselines on Financial PhraseBank & Agent-specific scores predict regional inflation \& employment better than ``hawkish/dovish'' \\ \hline
When to reach for it & Fast firm-level sentiment at scale & Central-bank communication studies \\ \hline
\end{tabular}
\end{center}

\medskip
\textbf{Both are fine-tuned BERTs} --- the architecture is from Topic 7.2; the pre-training/fine-tuning recipe is from Topic 7.3a. T4 cares only about \emph{using} them: load weights, feed sentences, get scores.
\end{frame}

\begin{frame}{Domain Models vs.\ Frontier LLMs: No Single Winner}
\textbf{Hansen \& Kazinnik (2023): ``GPT-4 reads Fedspeak.''} Zero-shot GPT-4 classifies FOMC sentence stance at expert-level accuracy; beats fine-tuned BERT/RoBERTa; provides human-consistent explanations.

\smallskip
\textbf{Gambacorta et al.\ (BIS WP 1215, 2024): CB-LMs.} Domain-adapted encoder beats generic BERT on FOMC stance; GPT-4 still wins on complex narrative tasks.

\smallskip
\textbf{Practical decision rule:}
\begin{itemize}\itemsep1pt
  \item \textbf{Sentence-level sentiment at scale, fixed taxonomy} $\to$ FinBERT / CentralBankRoBERTa (cheap, deterministic, citable).
  \item \textbf{Open-ended interpretation, narrative, novel category} $\to$ frontier LLM (GPT-4, Claude) with a logged prompt.
  \item \textbf{In doubt} $\to$ run both and report agreement (Cohen's $\kappa$).
\end{itemize}

\smallskip
\textit{No single model dominates --- the validation checklist (later in this topic) handles the choice.}
\end{frame}

\begin{frame}{Case 3b --- PaECTER for Patent Innovation}
\textbf{Ganguli, Lin \& Reynolds (2024), NBER WP}

\medskip
\begin{itemize}
\item \textbf{Goal:} measure innovation dynamics from patent text using transformer-based embeddings.

\medskip
\item \textbf{Model: PaECTER} (Patent Embeddings from Citation-Enhanced Transformer Representations)
\begin{itemize}
    \item Initialize from BERT-base
    \item Fine-tune on patent corpora
    \item Use citation pairs as a supervised signal during fine-tuning (related patents pulled together in embedding space)
\end{itemize}

\medskip
\item \textbf{Findings:}
\begin{itemize}
    \item Significantly more accurate similarity judgments than TF-IDF or LDA
    \item Distinguishes incremental from radical innovations
    \item Long-run trend: patents are increasingly diverging in semantic space (broader exploration)
\end{itemize}

\medskip
\item \textbf{Take-away:} fine-tuned transformer embeddings give a richer, citation-aware view of technological evolution than purely text-based or count-based methods.
\end{itemize}
\end{frame}

\begin{frame}{Frontier (2024--25): LLMs for Central Bank Communication}
\begin{itemize}
\item \textbf{CB-LMs} (Gambacorta et al., BIS WP 1215, 2024)
\begin{itemize}
    \item Domain-adapted encoder retrained on 37{,}000 central bank papers + 18{,}000 speeches
    \item Outperforms generic BERT on FOMC stance classification
    \item But: GPT-4 wins on complex narrative tasks --- no single model dominates
\end{itemize}

\medskip
\item \textbf{GPT-4 for Fedspeak} (Hansen \& Kazinnik, 2023)
\begin{itemize}
    \item GPT-4 classifies FOMC statement sentences at \emph{expert-level} accuracy
    \item Outperforms fine-tuned BERT/RoBERTa; provides human-consistent explanations
    \item Zero-shot performance collapses domain-adaptation barrier for many tasks
\end{itemize}

\medskip
\item \textbf{GPT-4o on 25{,}000 FOMC Sentences} (Federal Reserve FEDS Note, December 2024)
\begin{itemize}
    \item Topic classification across 2010--2024 minutes at scale
    \item Detects ``financial stability'' discourse spikes in 2020 and 2023 banking crises
    \item LLMs reveal \emph{thematic shifts} that sentiment scores miss
\end{itemize}

\medskip
\item \textbf{74{,}882 Central Bank Documents} (Silva, Moriya \& Veyrune, IMF WP 2025/109)
\begin{itemize}
    \item 169 central banks, 1884--2025; LLM-coded stance + directional communication index
    \item Forecasts future policy rate changes; documents shift to forward-looking language
\end{itemize}
\end{itemize}
\end{frame}

\begin{frame}{Frontier (2024--25): LLMs in Asset Pricing}
\begin{itemize}
\item \textbf{Expected Returns and LLMs} (Chen, Kelly \& Xiu, NBER WP 31502, 2024)
\begin{itemize}
    \item LLM embeddings of financial news $\Rightarrow$ cross-sectional equity return prediction
    \item Substantially outperforms TF-IDF and momentum; gains largest on \emph{negation} and complex narratives
    \item Evidence of slow price adjustment to information that only LLM embeddings extract
\end{itemize}

\medskip
\item \textbf{AI Asset Pricing Models} (Kelly, Kuznetsov, Malamud \& Xu, NBER WP 33351, 2025)
\begin{itemize}
    \item Transformer architecture embedded directly in the stochastic discount factor $M_{t+1}$
    \item Multi-head attention learns cross-asset dependencies from 132 stock characteristics
    \item Large reductions in pricing errors vs.\ previous ML factor models
\end{itemize}

\medskip
\item \textbf{Asset Embeddings} (Gabaix, Koijen, Richmond \& Yogo, NBER WP 33651, 2025)
\begin{itemize}
    \item Embeddings learned from \emph{investor portfolio holdings} via skip-gram (same math as Word2Vec)
    \item Arithmetic: \texttt{tech $+$ growth $-$ hardware $\approx$ software firms}
    \item GPT-4 generates economic narratives for firm clusters in the learned embedding space
\end{itemize}

\medskip
\item \textbf{Key insight:} the same transformer and embedding toolkit that encodes language can encode financial structure --- text and quantitative data are no longer separate worlds.
\end{itemize}
\end{frame}

\begin{frame}{The Application Zoo, Organized by Econometric Role}
\textbf{The same text can play three very different roles. Pick the representation to match the role.}
\medskip

\begin{center}
\renewcommand{\arraystretch}{1.25}
\scriptsize
\begin{tabular}{|p{3.2cm}|p{3.7cm}|p{3.5cm}|p{3.2cm}|}
\hline
\textbf{Paper} & \textbf{Signal (regressor $m(x_i)$)} & \textbf{Outcome (behavior $y$)} & \textbf{Instrument / shock} \\
\hline
Baker--Bloom--Davis (2016) & Newspaper keyword EPU $\to$ VAR & --- & EPU surprises \\
\hline
Hoberg \& Phillips (2016) & 10-K TF-IDF cosine $\to$ product-market similarity & --- & --- \\
\hline
Hassan et al.\ (2019) & Bigram firm-level political risk $\to$ investment regression & --- & Quarterly call shocks \\
\hline
Li et al.\ (2021) & W2V on earnings calls $\to$ 5 culture dimensions & --- & --- \\
\hline
Lopez-Lira \& Tang (2023) & ChatGPT zero-shot sentiment on headlines $\to$ next-day returns & --- & --- \\
\hline
Hansen--McMahon--Prat (2018) & LDA topics & Committee speech change after transparency reform & 1993 reform date \\
\hline
Hartley (2025) & LLM-reconstructed EPU (19th c., OCR) & --- & War-news shocks \\
\hline
\end{tabular}
\end{center}

\smallskip
\footnotesize
\begin{itemize}\itemsep0pt
\item \emph{Signal column dominates the literature} --- it's the easiest role and the loudest result.
\item \emph{Outcome} requires a behavior to be \emph{caused} by something else; harder but cleaner identification.
\item \emph{Instrument / shock} use needs the exclusion restriction to survive the prompt/fine-tuning recipe --- document it.
\item \textbf{Representation/role pairing.} Measurement (signal) $\to$ transparent (TF-IDF, dictionaries); behavior $\to$ contextual embeddings; instrument $\to$ frozen-snapshot LLM with archived prompt.
\end{itemize}
\end{frame}

\begin{frame}{Validation Caveats for Econometricians}
\begin{itemize}
\item \textbf{Cosine similarity is not innocent.} Three common confounds inflate similarity between firm/policy texts:
\begin{itemize}
    \item \textbf{Industry jargon:} firms in the same SIC code share boilerplate language
    \item \textbf{Templates:} regulatory filings re-use chunks of text verbatim
    \item \textbf{Length:} longer documents have smoother averages and look more similar
\end{itemize}

\medskip
\item \textbf{Standard tools for robustness:}
\begin{itemize}
    \item Permutation tests --- shuffle labels within an industry, check if observed similarity exceeds the within-industry baseline
    \item Residual embeddings --- subtract the industry mean before measuring similarity
    \item Bootstrap the entire pipeline (text $\rightarrow$ embedding $\rightarrow$ measurement $\rightarrow$ regression) for honest standard errors
\end{itemize}

\medskip
\item \textbf{Dictionary methods are not innocent either.} The Loughran-McDonald (2011) finance-specific dictionary corrects mismatches in the older Harvard IV-4 list (e.g.\ ``liability'' is not negative in financial accounting).

\medskip
\item \textbf{Always validate} a sample of model outputs against human coding before scaling up.
\end{itemize}
\end{frame}

\begin{frame}{A Six-Step Workflow for LLM-Assisted Measurement}
\begin{enumerate}
\item \textbf{Pre-alignment.} Hand-code 100--200 examples; iterate on the prompt or fine-tuning recipe until you reach acceptable agreement.
\medskip
\item \textbf{Full annotation.} Run the model on the entire sample with a frozen prompt and a frozen model version.
\medskip
\item \textbf{Gold-standard audit.} Blind-code 10\% of the sample and compute Cohen's $\kappa$ (publication threshold $\kappa \ge 0.75$).
\medskip
\item \textbf{Category-level reporting.} Don't just report overall accuracy; show the confusion matrix by class. Minority classes are often where models fail silently.
\medskip
\item \textbf{External validity check.} Link the text-derived measure to a real economic outcome (returns, investment, employment) before treating it as a measurement.
\medskip
\item \textbf{Disclosure.} Report prompt, model version, temperature ($\tau = 0$), date, all preprocessing, and provide code. Reproducibility is essential when the model behind the API can change without notice.
\end{enumerate}
\end{frame}

\begin{frame}{Worked Example: Applying the Checklist to FOMC Sentiment}
\textbf{Imagined paper:} ``GPT-4 sentiment on FOMC statements predicts 10-year yields.''

\medskip
\begin{center}
\renewcommand{\arraystretch}{1.2}
\footnotesize
\begin{tabular}{|p{2.4cm}|p{4.0cm}|p{6.6cm}|}
\hline
\textbf{Step} & \textbf{What you do} & \textbf{What goes in the appendix} \\ \hline
1. Pre-align & Hand-code 200 sentences hawk/dove/neutral & Prompt v3 (frozen); $\kappa$ vs.\ author = 0.81 \\ \hline
2. Annotate & GPT-4-2024-06-13, $\tau=0$, full corpus & Cost, runtime, retries, sha256 of inputs \\ \hline
3. Gold audit & Blind-code 10\% by RA & $\kappa = 0.78$ overall; 0.62 on ``neutral'' \\ \hline
4. By-class & Confusion matrix & Hawkish recall 0.91; neutral recall 0.66 (flagged) \\ \hline
5. External validity & Predict $\Delta y_{10}$ at FOMC dates & $\hat\beta = 4.2$ bp per unit, $t = 3.1$ \\ \hline
6. Disclose & Prompt + code on Zenodo & DOI, model card, replication run-time \\ \hline
\end{tabular}
\end{center}

\medskip
\textbf{LLM-as-judge connection.} A second LLM as a \texttt{cross-checker} sub-agent: same input, different prompt, flag disagreements for human review --- the same construction-to-coefficient logic used in AI-exposure scoring (Topic 7.3b).
\end{frame}

\begin{frame}{Disclosure Norms for AI-Assisted Empirical Work}
\footnotesize
\textbf{Minimum reproducibility package.}
\begin{itemize}\itemsep1pt
  \item \textbf{Model identity.} Provider, model name, exact version string, call date.
  \item \textbf{Prompt(s).} Frozen text in a versioned file; not paraphrased in the paper.
  \item \textbf{Sampling parameters.} Temperature, top-$p$, max tokens, seed if available.
  \item \textbf{Inputs.} sha256 of every input file; document any preprocessing.
  \item \textbf{Outputs.} Raw model outputs preserved \emph{before} any post-processing.
  \item \textbf{Code.} End-to-end script with retry logic and the exact API client version.
\end{itemize}

\smallskip
\textbf{AEA policy.} The October 2024 Data and Code Availability update explicitly covers AI-assisted analysis: report when an LLM was used, for what step, and provide artifacts that let a reader reproduce the result given the same model.

\smallskip
\textbf{IRB and proprietary text.} Survey transcripts, internal memos, and clinical text often cannot be sent to a hosted API. Run an open-weights model locally, or obtain IRB sign-off on the provider's data-handling terms.

\smallskip
\textit{The closed-source LLM is the most-changed equipment in your empirical paper. Document it like the dataset it is.}
\end{frame}

\begin{frame}{Topic 7.4 Wrap-Up --- and the Bridge to Lec08/Lec09/Lec10}
\begin{itemize}
\item \textbf{Three econometric roles of text.} Signal (regressor), outcome (behavior), instrument (shock). The role determines the representation: transparent for measurement, contextual for behavior, frozen-snapshot for instruments.

\medskip
\item \textbf{Six-step measurement workflow.} Pre-align $\to$ full annotation $\to$ gold-standard audit $\to$ category-level reporting $\to$ external validity $\to$ disclosure. Cohen's $\kappa \ge 0.75$ is the publication bar.

\medskip
\item \textbf{Disclosure norms (AEA, October 2024).} Model identity, prompt, sampling parameters, input hashes, raw outputs, end-to-end code. Treat the LLM call as you would a survey instrument.

\medskip
\item \textbf{What's next.}
\begin{itemize}
    \item \textbf{Lec08 (RAG)} --- ground the LLM in your private corpus. Solves T3b's hallucination + cutoff + context-window limits via retrieval; the token-vs-doc-embedding frame from T2b is the geometric prerequisite.
    \item \textbf{Lec09 (Agentic AI)} --- ground the LLM in actions. JSON schemas from T3a become tool calls; ICL-as-Bayesian (T3a) returns as the context dilemma.
    \item \textbf{Lec10 (Case Studies)} --- end-to-end empirical papers built with the Lec07--Lec09 toolkit.
\end{itemize}

\medskip
\centering\emph{Language engine (Lec07) $+$ grounded memory (Lec08) $+$ grounded action (Lec09) $=$ auditable economic measurement.}
\end{itemize}
\end{frame}

\end{document}
